\section{Tűzfalak és forgalom-szabályozás}

A host-alapú tűzfalak (firewalld, nftables) és a szabályzatok tervezése fontos része a védelemnek.

\subsection{firewalld alapok}
\begin{lstlisting}[language=bash]
# Check status
firewall-cmd --state
# Allow a service at runtime
firewall-cmd --zone=public --add-service=http
# Persist the change
firewall-cmd --permanent --zone=public --add-service=http
firewall-cmd --runtime-to-permanent
# List zone settings
firewall-cmd --list-all --zone=public
\end{lstlisting}

\subsection{nftables és netfilter}
\begin{lstlisting}[language=bash]
# Full ruleset
nft list ruleset
# Create a new table and chain
nft add table ip filter
nft add chain ip filter input { type filter hook input priority 0; }
\end{lstlisting}

Tervezési szempontok:
\begin{itemize}
  \item Kizárólag szükséges szolgáltatásokat engedélyezzünk.
  \item Zóna alapú megközelítés a host szerepének megfelelően (dmz, internal, trusted).
  \item Logoljuk a drop-olt csomagokat, de figyeljünk az adatmennyiségre.
\end{itemize}
